\documentclass[8pt,a4paper]{extarticle}
% ============================================================
% Gemeinsamer Preamble fuer alle 5 Open-Book-Spickzettel.
% Wird aus jedem klausur.tex via % ============================================================
% Gemeinsamer Preamble fuer alle 5 Open-Book-Spickzettel.
% Wird aus jedem klausur.tex via % ============================================================
% Gemeinsamer Preamble fuer alle 5 Open-Book-Spickzettel.
% Wird aus jedem klausur.tex via \input{../preamble.tex} geladen.
% Layout: A4 hochkant, 3 Spalten, 8pt, sehr kompakt.
% ============================================================

\usepackage[a4paper,margin=8mm,top=7mm,bottom=7mm,includefoot]{geometry}
\usepackage[utf8]{inputenc}
\usepackage[T1]{fontenc}
\usepackage[ngerman]{babel}
\usepackage{lmodern}
\usepackage{microtype}

\usepackage{amsmath,amssymb,amsfonts,mathtools}
\usepackage{multicol}
\usepackage{enumitem}
\usepackage{array,tabularx,booktabs,makecell}
\usepackage[dvipsnames,table]{xcolor}
\usepackage[explicit]{titlesec}
\usepackage{etoolbox}
\usepackage{fancyhdr}
\usepackage{lastpage}

% ---------- Farben ----------
\definecolor{secblue}{HTML}{1f4e79}
\definecolor{subblue}{HTML}{2e75b6}
\definecolor{boxbg}{HTML}{f2f2f2}
\definecolor{accent}{HTML}{c00000}
\definecolor{rulegray}{HTML}{888888}

% ---------- Spalten ----------
\setlength{\columnsep}{3mm}
\setlength{\columnseprule}{0.3pt}
\def\columnseprulecolor{\color{rulegray}}

% ---------- Listen kompakt ----------
\setlist{nosep,leftmargin=1.1em,labelsep=0.35em,topsep=0.2ex,partopsep=0pt}
\setlist[itemize]{label=\textbullet}

% ---------- Abstaende ----------
\setlength{\parindent}{0pt}
\setlength{\parskip}{0.15ex}
\setlength{\abovedisplayskip}{0.3ex}
\setlength{\belowdisplayskip}{0.3ex}
\setlength{\abovedisplayshortskip}{0.2ex}
\setlength{\belowdisplayshortskip}{0.2ex}
\setlength{\jot}{1pt}

% ---------- Tabellen kompakt ----------
\renewcommand{\arraystretch}{1.05}
\setlength{\tabcolsep}{3pt}
\setlength{\arrayrulewidth}{0.3pt}
\arrayrulecolor{rulegray}

% ---------- Section-Stil (kein Vollblock; dick + farbige Linie darunter) ----------
\titleformat{\section}
  {\large\bfseries\color{secblue}}
  {}{0pt}{#1}
  [\vspace{-0.4ex}{\color{secblue}\hrule height 0.8pt}\vspace{0.1ex}]
\titlespacing*{\section}{0pt}{1.4ex plus 0.4ex}{0.5ex}

\titleformat{\subsection}
  {\small\bfseries\color{subblue}}
  {}{0pt}{#1}
\titlespacing*{\subsection}{0pt}{0.9ex plus 0.2ex}{0.2ex}

\titleformat{\subsubsection}
  {\small\bfseries}
  {}{0pt}{#1}
\titlespacing*{\subsubsection}{0pt}{0.5ex}{0.1ex}

% ---------- Mini-Helfer ----------

% \fact{...} Hervorgehobene Regel (gerahmt, hellgrau)
\newcommand{\fact}[1]{%
  \begingroup\setlength{\fboxsep}{2pt}%
  \colorbox{boxbg}{\parbox{\dimexpr\linewidth-2\fboxsep}{\small #1\strut}}%
  \endgroup\par\vspace{0.2ex}}

% \warn{...} Stolperfalle (roter Strich links)
\newcommand{\warn}[1]{%
  \begingroup\setlength{\fboxsep}{2pt}%
  \noindent{\color{accent}\rule{2pt}{\baselineskip}}\hspace{2pt}%
  \begin{minipage}[t]{\dimexpr\linewidth-6pt}\small #1\end{minipage}%
  \endgroup\par\vspace{0.2ex}}

% Math-Shortcuts
\newcommand{\R}{\mathbb{R}}
\newcommand{\N}{\mathbb{N}}
\newcommand{\Z}{\mathbb{Z}}
\newcommand{\Q}{\mathbb{Q}}
\newcommand{\C}{\mathbb{C}}
\DeclareMathOperator{\rg}{rg}
\DeclareMathOperator{\spur}{tr}
\DeclareMathOperator{\im}{im}
\DeclareMathOperator{\diag}{diag}
\DeclareMathOperator{\sign}{sgn}
\DeclareMathOperator*{\argmin}{arg\,min}
\DeclareMathOperator*{\argmax}{arg\,max}
\newcommand{\trans}{^{\!\top}}
\newcommand{\norm}[1]{\left\lVert #1 \right\rVert}
\newcommand{\abs}[1]{\left\lvert #1 \right\rvert}
\newcommand{\inner}[2]{\langle #1,\,#2\rangle}
\newcommand{\dx}{\,\mathrm{d}x}
\newcommand{\dy}{\,\mathrm{d}y}
\newcommand{\dt}{\,\mathrm{d}t}

% Zeigt eine Display-Formel zentriert; skaliert nur runter wenn zu breit.
\newlength{\fitw}
\newcommand{\fitline}[1]{%
  \par\nointerlineskip
  \settowidth{\fitw}{$\displaystyle #1$}%
  \ifdim\fitw>\linewidth
    \noindent\resizebox{\linewidth}{!}{$\displaystyle #1$}\par
  \else
    \[\displaystyle #1\]%
  \fi}

% Tabular fuer fixe Spaltenbreiten (Spickzettel-typisch)
\newenvironment{tightt}[1]
  {\par\noindent\begin{tabular}{#1}}
  {\end{tabular}\par}
% Fuer X-Spalten direkt \begin{tabularx}{\linewidth}{...} verwenden
% (kann nicht in \newenvironment gewrappt werden -- TX@get@body Konflikt)

% Globale Schutzschalter gegen Overflow
\emergencystretch=2em
\tolerance=2000
\hbadness=10000
\hfuzz=2pt
\vfuzz=2pt

% ---------- Kopf-/Fusszeile ----------
\pagestyle{fancy}
\fancyhf{}
\renewcommand{\headrulewidth}{0pt}
\renewcommand{\footrulewidth}{0pt}
\fancyfoot[L]{\scriptsize\color{rulegray}\textit{\sheettitle}}
\fancyfoot[R]{\scriptsize\color{rulegray}Seite \thepage\,/\,\pageref{LastPage}}
\setlength{\footskip}{8pt}

% Titel-Header oben auf Seite 1
\newcommand{\sheetheader}[2]{%
  {\Large\bfseries\color{secblue} #1}\hfill
  {\small\color{rulegray} #2}\par
  \vspace{0.4ex}{\color{secblue}\hrule height 0.7pt}\vspace{1ex}}

% Default-Titel (wird ueberschrieben)
\newcommand{\sheettitle}{Open-Book Spickzettel}
 geladen.
% Layout: A4 hochkant, 3 Spalten, 8pt, sehr kompakt.
% ============================================================

\usepackage[a4paper,margin=8mm,top=7mm,bottom=7mm,includefoot]{geometry}
\usepackage[utf8]{inputenc}
\usepackage[T1]{fontenc}
\usepackage[ngerman]{babel}
\usepackage{lmodern}
\usepackage{microtype}

\usepackage{amsmath,amssymb,amsfonts,mathtools}
\usepackage{multicol}
\usepackage{enumitem}
\usepackage{array,tabularx,booktabs,makecell}
\usepackage[dvipsnames,table]{xcolor}
\usepackage[explicit]{titlesec}
\usepackage{etoolbox}
\usepackage{fancyhdr}
\usepackage{lastpage}

% ---------- Farben ----------
\definecolor{secblue}{HTML}{1f4e79}
\definecolor{subblue}{HTML}{2e75b6}
\definecolor{boxbg}{HTML}{f2f2f2}
\definecolor{accent}{HTML}{c00000}
\definecolor{rulegray}{HTML}{888888}

% ---------- Spalten ----------
\setlength{\columnsep}{3mm}
\setlength{\columnseprule}{0.3pt}
\def\columnseprulecolor{\color{rulegray}}

% ---------- Listen kompakt ----------
\setlist{nosep,leftmargin=1.1em,labelsep=0.35em,topsep=0.2ex,partopsep=0pt}
\setlist[itemize]{label=\textbullet}

% ---------- Abstaende ----------
\setlength{\parindent}{0pt}
\setlength{\parskip}{0.15ex}
\setlength{\abovedisplayskip}{0.3ex}
\setlength{\belowdisplayskip}{0.3ex}
\setlength{\abovedisplayshortskip}{0.2ex}
\setlength{\belowdisplayshortskip}{0.2ex}
\setlength{\jot}{1pt}

% ---------- Tabellen kompakt ----------
\renewcommand{\arraystretch}{1.05}
\setlength{\tabcolsep}{3pt}
\setlength{\arrayrulewidth}{0.3pt}
\arrayrulecolor{rulegray}

% ---------- Section-Stil (kein Vollblock; dick + farbige Linie darunter) ----------
\titleformat{\section}
  {\large\bfseries\color{secblue}}
  {}{0pt}{#1}
  [\vspace{-0.4ex}{\color{secblue}\hrule height 0.8pt}\vspace{0.1ex}]
\titlespacing*{\section}{0pt}{1.4ex plus 0.4ex}{0.5ex}

\titleformat{\subsection}
  {\small\bfseries\color{subblue}}
  {}{0pt}{#1}
\titlespacing*{\subsection}{0pt}{0.9ex plus 0.2ex}{0.2ex}

\titleformat{\subsubsection}
  {\small\bfseries}
  {}{0pt}{#1}
\titlespacing*{\subsubsection}{0pt}{0.5ex}{0.1ex}

% ---------- Mini-Helfer ----------

% \fact{...} Hervorgehobene Regel (gerahmt, hellgrau)
\newcommand{\fact}[1]{%
  \begingroup\setlength{\fboxsep}{2pt}%
  \colorbox{boxbg}{\parbox{\dimexpr\linewidth-2\fboxsep}{\small #1\strut}}%
  \endgroup\par\vspace{0.2ex}}

% \warn{...} Stolperfalle (roter Strich links)
\newcommand{\warn}[1]{%
  \begingroup\setlength{\fboxsep}{2pt}%
  \noindent{\color{accent}\rule{2pt}{\baselineskip}}\hspace{2pt}%
  \begin{minipage}[t]{\dimexpr\linewidth-6pt}\small #1\end{minipage}%
  \endgroup\par\vspace{0.2ex}}

% Math-Shortcuts
\newcommand{\R}{\mathbb{R}}
\newcommand{\N}{\mathbb{N}}
\newcommand{\Z}{\mathbb{Z}}
\newcommand{\Q}{\mathbb{Q}}
\newcommand{\C}{\mathbb{C}}
\DeclareMathOperator{\rg}{rg}
\DeclareMathOperator{\spur}{tr}
\DeclareMathOperator{\im}{im}
\DeclareMathOperator{\diag}{diag}
\DeclareMathOperator{\sign}{sgn}
\DeclareMathOperator*{\argmin}{arg\,min}
\DeclareMathOperator*{\argmax}{arg\,max}
\newcommand{\trans}{^{\!\top}}
\newcommand{\norm}[1]{\left\lVert #1 \right\rVert}
\newcommand{\abs}[1]{\left\lvert #1 \right\rvert}
\newcommand{\inner}[2]{\langle #1,\,#2\rangle}
\newcommand{\dx}{\,\mathrm{d}x}
\newcommand{\dy}{\,\mathrm{d}y}
\newcommand{\dt}{\,\mathrm{d}t}

% Zeigt eine Display-Formel zentriert; skaliert nur runter wenn zu breit.
\newlength{\fitw}
\newcommand{\fitline}[1]{%
  \par\nointerlineskip
  \settowidth{\fitw}{$\displaystyle #1$}%
  \ifdim\fitw>\linewidth
    \noindent\resizebox{\linewidth}{!}{$\displaystyle #1$}\par
  \else
    \[\displaystyle #1\]%
  \fi}

% Tabular fuer fixe Spaltenbreiten (Spickzettel-typisch)
\newenvironment{tightt}[1]
  {\par\noindent\begin{tabular}{#1}}
  {\end{tabular}\par}
% Fuer X-Spalten direkt \begin{tabularx}{\linewidth}{...} verwenden
% (kann nicht in \newenvironment gewrappt werden -- TX@get@body Konflikt)

% Globale Schutzschalter gegen Overflow
\emergencystretch=2em
\tolerance=2000
\hbadness=10000
\hfuzz=2pt
\vfuzz=2pt

% ---------- Kopf-/Fusszeile ----------
\pagestyle{fancy}
\fancyhf{}
\renewcommand{\headrulewidth}{0pt}
\renewcommand{\footrulewidth}{0pt}
\fancyfoot[L]{\scriptsize\color{rulegray}\textit{\sheettitle}}
\fancyfoot[R]{\scriptsize\color{rulegray}Seite \thepage\,/\,\pageref{LastPage}}
\setlength{\footskip}{8pt}

% Titel-Header oben auf Seite 1
\newcommand{\sheetheader}[2]{%
  {\Large\bfseries\color{secblue} #1}\hfill
  {\small\color{rulegray} #2}\par
  \vspace{0.4ex}{\color{secblue}\hrule height 0.7pt}\vspace{1ex}}

% Default-Titel (wird ueberschrieben)
\newcommand{\sheettitle}{Open-Book Spickzettel}
 geladen.
% Layout: A4 hochkant, 3 Spalten, 8pt, sehr kompakt.
% ============================================================

\usepackage[a4paper,margin=8mm,top=7mm,bottom=7mm,includefoot]{geometry}
\usepackage[utf8]{inputenc}
\usepackage[T1]{fontenc}
\usepackage[ngerman]{babel}
\usepackage{lmodern}
\usepackage{microtype}

\usepackage{amsmath,amssymb,amsfonts,mathtools}
\usepackage{multicol}
\usepackage{enumitem}
\usepackage{array,tabularx,booktabs,makecell}
\usepackage[dvipsnames,table]{xcolor}
\usepackage[explicit]{titlesec}
\usepackage{etoolbox}
\usepackage{fancyhdr}
\usepackage{lastpage}

% ---------- Farben ----------
\definecolor{secblue}{HTML}{1f4e79}
\definecolor{subblue}{HTML}{2e75b6}
\definecolor{boxbg}{HTML}{f2f2f2}
\definecolor{accent}{HTML}{c00000}
\definecolor{rulegray}{HTML}{888888}

% ---------- Spalten ----------
\setlength{\columnsep}{3mm}
\setlength{\columnseprule}{0.3pt}
\def\columnseprulecolor{\color{rulegray}}

% ---------- Listen kompakt ----------
\setlist{nosep,leftmargin=1.1em,labelsep=0.35em,topsep=0.2ex,partopsep=0pt}
\setlist[itemize]{label=\textbullet}

% ---------- Abstaende ----------
\setlength{\parindent}{0pt}
\setlength{\parskip}{0.15ex}
\setlength{\abovedisplayskip}{0.3ex}
\setlength{\belowdisplayskip}{0.3ex}
\setlength{\abovedisplayshortskip}{0.2ex}
\setlength{\belowdisplayshortskip}{0.2ex}
\setlength{\jot}{1pt}

% ---------- Tabellen kompakt ----------
\renewcommand{\arraystretch}{1.05}
\setlength{\tabcolsep}{3pt}
\setlength{\arrayrulewidth}{0.3pt}
\arrayrulecolor{rulegray}

% ---------- Section-Stil (kein Vollblock; dick + farbige Linie darunter) ----------
\titleformat{\section}
  {\large\bfseries\color{secblue}}
  {}{0pt}{#1}
  [\vspace{-0.4ex}{\color{secblue}\hrule height 0.8pt}\vspace{0.1ex}]
\titlespacing*{\section}{0pt}{1.4ex plus 0.4ex}{0.5ex}

\titleformat{\subsection}
  {\small\bfseries\color{subblue}}
  {}{0pt}{#1}
\titlespacing*{\subsection}{0pt}{0.9ex plus 0.2ex}{0.2ex}

\titleformat{\subsubsection}
  {\small\bfseries}
  {}{0pt}{#1}
\titlespacing*{\subsubsection}{0pt}{0.5ex}{0.1ex}

% ---------- Mini-Helfer ----------

% \fact{...} Hervorgehobene Regel (gerahmt, hellgrau)
\newcommand{\fact}[1]{%
  \begingroup\setlength{\fboxsep}{2pt}%
  \colorbox{boxbg}{\parbox{\dimexpr\linewidth-2\fboxsep}{\small #1\strut}}%
  \endgroup\par\vspace{0.2ex}}

% \warn{...} Stolperfalle (roter Strich links)
\newcommand{\warn}[1]{%
  \begingroup\setlength{\fboxsep}{2pt}%
  \noindent{\color{accent}\rule{2pt}{\baselineskip}}\hspace{2pt}%
  \begin{minipage}[t]{\dimexpr\linewidth-6pt}\small #1\end{minipage}%
  \endgroup\par\vspace{0.2ex}}

% Math-Shortcuts
\newcommand{\R}{\mathbb{R}}
\newcommand{\N}{\mathbb{N}}
\newcommand{\Z}{\mathbb{Z}}
\newcommand{\Q}{\mathbb{Q}}
\newcommand{\C}{\mathbb{C}}
\DeclareMathOperator{\rg}{rg}
\DeclareMathOperator{\spur}{tr}
\DeclareMathOperator{\im}{im}
\DeclareMathOperator{\diag}{diag}
\DeclareMathOperator{\sign}{sgn}
\DeclareMathOperator*{\argmin}{arg\,min}
\DeclareMathOperator*{\argmax}{arg\,max}
\newcommand{\trans}{^{\!\top}}
\newcommand{\norm}[1]{\left\lVert #1 \right\rVert}
\newcommand{\abs}[1]{\left\lvert #1 \right\rvert}
\newcommand{\inner}[2]{\langle #1,\,#2\rangle}
\newcommand{\dx}{\,\mathrm{d}x}
\newcommand{\dy}{\,\mathrm{d}y}
\newcommand{\dt}{\,\mathrm{d}t}

% Zeigt eine Display-Formel zentriert; skaliert nur runter wenn zu breit.
\newlength{\fitw}
\newcommand{\fitline}[1]{%
  \par\nointerlineskip
  \settowidth{\fitw}{$\displaystyle #1$}%
  \ifdim\fitw>\linewidth
    \noindent\resizebox{\linewidth}{!}{$\displaystyle #1$}\par
  \else
    \[\displaystyle #1\]%
  \fi}

% Tabular fuer fixe Spaltenbreiten (Spickzettel-typisch)
\newenvironment{tightt}[1]
  {\par\noindent\begin{tabular}{#1}}
  {\end{tabular}\par}
% Fuer X-Spalten direkt \begin{tabularx}{\linewidth}{...} verwenden
% (kann nicht in \newenvironment gewrappt werden -- TX@get@body Konflikt)

% Globale Schutzschalter gegen Overflow
\emergencystretch=2em
\tolerance=2000
\hbadness=10000
\hfuzz=2pt
\vfuzz=2pt

% ---------- Kopf-/Fusszeile ----------
\pagestyle{fancy}
\fancyhf{}
\renewcommand{\headrulewidth}{0pt}
\renewcommand{\footrulewidth}{0pt}
\fancyfoot[L]{\scriptsize\color{rulegray}\textit{\sheettitle}}
\fancyfoot[R]{\scriptsize\color{rulegray}Seite \thepage\,/\,\pageref{LastPage}}
\setlength{\footskip}{8pt}

% Titel-Header oben auf Seite 1
\newcommand{\sheetheader}[2]{%
  {\Large\bfseries\color{secblue} #1}\hfill
  {\small\color{rulegray} #2}\par
  \vspace{0.4ex}{\color{secblue}\hrule height 0.7pt}\vspace{1ex}}

% Default-Titel (wird ueberschrieben)
\newcommand{\sheettitle}{Open-Book Spickzettel}

\renewcommand{\sheettitle}{Numerische Methoden}

\begin{document}
\sheetheader{Numerische Methoden}{Open-Book Formelsammlung -- DHBW WDS125}

\begin{multicols*}{3}

% =================================================
\section{Fehler \& Kondition}
% =================================================
\subsection{Maschinenzahlen (IEEE 754)}
Double: $\pm 1.m\cdot 2^e$, 52 Mant.-, 11 Exp.-Bits.\\
$\epsilon_{\text{mach}}=2^{-52}\!\approx\!2{,}22\cdot 10^{-16}$.\\
Spez.: $\pm 0,\ \pm\infty$, NaN, Denormal.

\subsection{Fehlerarten}
\begin{tabularx}{\linewidth}{@{}l@{\ }X@{}}
Datenfehler & Eingabeungenauigk. \\
Modellfehler & Physik vereinfacht \\
Diskretisier.\ & Verfahren approx. \\
Rundungsfehler & endliche Pr\"azision \\
Abbruchfehler & Iteration zu fr\"uh \\
\end{tabularx}
Absolut $|\tilde x-x|$;\ relativ $|\tilde x-x|/|x|$.

\subsection{Ausl\"oschung}
$a-b$ bei $a\!\approx\!b$ vernichtet Stellen. Umformen:
\begin{itemize}
\item $1-\cos x = 2\sin^2(x/2)$
\item $\sqrt{a+h}-\sqrt a = h/(\sqrt{a+h}+\sqrt a)$
\item $e^x-1$: \texttt{expm1(x)} bzw.\ Taylor
\end{itemize}

\subsection{Konditionszahl}
\begin{itemize}
\item Skalar: $\kappa_{\text{rel}}(f,x)=|x f'(x)/f(x)|$
\item Matrix: $\kappa(A)=\|A\|\cdot\|A^{-1}\|$,\ $\kappa_2=\sigma_1/\sigma_n$
\item $\kappa\!\approx\!10^{16}$ in double $\Rightarrow$ keine signif.\ Stelle mehr
\end{itemize}
\fact{Problem hat \textbf{Kondition}, Verfahren hat \textbf{Stabilit\"at}. Bsp: Gauss ohne Pivot bei kleinem Diag.\ instabil.}

% =================================================
\section{Lineare Systeme}
% =================================================
\subsection{Direkt}
Gauss mit Spaltenpivot: $PA=LU$, $\tfrac23 n^3$.\\
Cholesky (SPD): $\tfrac13 n^3$.\\
QR: Least Squares.

\subsection{Iterativ stationaer}
Splitting $A=M-N$: $x_{k+1}=M^{-1}(Nx_k+b)$.\\
Konv.\ $\Leftrightarrow$ Spektralradius $\rho(M^{-1}N)<1$.\\[0.3ex]
\textbf{Jacobi} ($M=D$):
\fitline{x_i^{(k+1)}=\tfrac1{a_{ii}}\!\left(b_i-\!\!\sum_{j\ne i}\!a_{ij} x_j^{(k)}\right)}
\textbf{Gauss-Seidel} (aktualisierte Komp.):
\fitline{x_i^{(k+1)}=\tfrac1{a_{ii}}\!\left(b_i-\!\sum_{j<i}\!a_{ij} x_j^{(k+1)}-\!\sum_{j>i}\!a_{ij} x_j^{(k)}\right)}
\textbf{SOR}: $x^{(k+1)}=(1{-}\omega)x^{(k)}+\omega\,x^{\text{GS},(k+1)}$, $\omega\!\in\!(0,2)$.\\
\fact{Diagonaldominanz $|a_{ii}|>\sum_{j\ne i}|a_{ij}|$ ist hinreichend f\"ur Jacobi/GS-Konvergenz.}

\subsection{CG (SPD) und Krylov}
\begin{itemize}
\item $r_0=b-Ax_0,\ p_0=r_0$
\item $\alpha_k=\frac{r_k\trans r_k}{p_k\trans A p_k}$
\item $x_{k+1}=x_k+\alpha_k p_k$
\item $r_{k+1}=r_k-\alpha_k A p_k$
\item $\beta_k=\frac{r_{k+1}\trans r_{k+1}}{r_k\trans r_k}$
\item $p_{k+1}=r_{k+1}+\beta_k p_k$
\end{itemize}
$\le n$ Schritte exakt; \textbf{GMRES} f\"ur allg.\ $A$ (Restart $m$).\\
\textbf{Stop}: $\|r^{(k)}\|/\|b\|<\tau$ (Residuum bevorzugt).

% =================================================
\section{Nichtlin.\ Gleichungen}
% =================================================
\subsection{Bisektion}
$f$ stetig, $f(a)f(b)<0$. Linear (Ord.\ 1) -- ca.\ 3.3 Iter./Dezimalstelle. Garantiert konv.

\subsection{Newton}
\fitline{x_{k+1}=x_k-f(x_k)/f'(x_k)}
Quadrat.\ lokal: $|e_{k+1}|\!\approx\!c|e_k|^2$.\\
Mehrdim.: $x_{k+1}=x_k-J_F(x_k)^{-1}F(x_k)$.\\
Heron $\sqrt 2$: $x_{k+1}=(x_k+2/x_k)/2$.

\subsection{Sekante}
\fitline{x_{k+1}=x_k-f(x_k)\cdot\tfrac{x_k-x_{k-1}}{f(x_k)-f(x_{k-1})}}
Ordnung $p=(1{+}\sqrt 5)/2\approx 1{,}618$.

\subsection{Fixpunkt}
$x=g(x)$, $x_{k+1}=g(x_k)$.\\
\textbf{Banach}: $g$ Kontraktion $|g(x){-}g(y)|\!\le\!L|x{-}y|$, $L\!<\!1$ ($\Leftrightarrow|g'|\!\le\!L\!<\!1$) $\Rightarrow$ eindeutiger Fixpunkt, Rate $L$.\ Quadrat.\ falls $g'(x^\star){=}0$.

\subsection{Konvergenzordnungen}
\begin{tabularx}{\linewidth}{@{}l@{\ }X@{}}
Bisektion & 1 (linear) \\
Sekante & $\approx 1{,}618$ \\
Newton & 2 (quadr.) \\
Fixpunkt & linear / quadr. \\
\end{tabularx}

% =================================================
\section{Interpolation}
% =================================================
Eindeutiges Polynom Grad $\le n$ zu $n{+}1$ St\"utzstellen.

\subsection{Lagrange}
\fitline{L_i(x)=\!\!\prod_{j\ne i}\!\tfrac{x-x_j}{x_i-x_j},\quad p(x)=\sum_i y_i L_i(x)}
$L_i(x_j)=\delta_{ij}$.

\subsection{Newton (div.\ Diff.)}
\fitline{p(x)=f[x_0]+f[x_0,x_1](x{-}x_0)+\dots}
\fitline{f[x_i,\dots,x_{i+k}]=\tfrac{f[x_{i+1}..]-f[x_i..]}{x_{i+k}-x_i}}

\subsection{Runge-Ph\"anomen}
Bei \"a\-qui\-di\-stan\-ten Knoten w\"achst der Fehler an den R\"andern.
\fact{Abhilfe: Tschebyscheff-Knoten $x_k=\cos\!\bigl(\tfrac{(2k+1)\pi}{2n+2}\bigr)$ oder Splines.}

\subsection{Splines}
\begin{itemize}
\item Linear $C^0$
\item Kubisch $C^2$ (Standard): nat\"urlich ($s''$=0 am Rand), clamped ($s'$ vorgeg.), period.\ ($s,s',s''$ gleich)
\end{itemize}

% =================================================
\section{Numerische Differentiation}
% =================================================
\begin{tabularx}{\linewidth}{@{}l@{\ }c@{}}
Vorw.\ $(f(x{+}h){-}f(x))/h$ & $O(h)$ \\
R\"uckw.\ $(f(x){-}f(x{-}h))/h$ & $O(h)$ \\
Zentral $(f(x{+}h){-}f(x{-}h))/(2h)$ & $O(h^2)$ \\
$f''$: $(f(x{+}h){-}2f(x){+}f(x{-}h))/h^2$ & $O(h^2)$ \\
\end{tabularx}
\textbf{Optimales $h$}: Vorw.\ $\sim\!\sqrt\epsilon\!\approx\!10^{-8}$; Zentral $\sim\!\sqrt[3]\epsilon\!\approx\!10^{-5}$.\\
\textbf{Richardson} ($O(h^2)$): $D_{\text{verb}}=D(h/2)+(D(h/2)-D(h))/3$.

% =================================================
\section{Quadratur}
% =================================================
\subsection{Newton-Cotes}
\textbf{Trapez}:
\fitline{\textstyle\int_a^b f\approx \tfrac{b-a}{2}\bigl(f(a)+f(b)\bigr)}
Fehler $-\tfrac{(b-a)^3}{12}f''(\xi)$, exakt Grad 1.\\[0.3ex]
\textbf{Simpson}:
\fitline{\textstyle\int_a^b f\approx \tfrac{b-a}{6}\bigl(f(a)+4 f(\tfrac{a+b}{2})+f(b)\bigr)}
Fehler $-\tfrac{(b-a)^5}{2880}f^{(4)}(\xi)$, exakt Grad 3.

\subsection{Zusammengesetzt ($h{=}(b{-}a)/n$)}
\fitline{T_n=\tfrac{h}{2}\bigl(f(a)+2\textstyle\sum_{i=1}^{n-1}\!f(x_i)+f(b)\bigr)\ \ O(h^2)}
\fitline{S_n=\tfrac{h}{3}\bigl(f(a)+4\textstyle\sum_{\text{u}}\!+\!2\textstyle\sum_{\text{g}}\!+f(b)\bigr)\ \ O(h^4)}
``u'' = ungerade Indizes, ``g'' = gerade.

\subsection{Gauss-Quadratur}
$\int_{-1}^1 f\approx\sum_{i=1}^n w_i f(x_i)$ mit $x_i$ = Nullst.\ Legendre; exakt f\"ur Grad $\le 2n{-}1$.\\
$[a,b]$: $x=\tfrac{b-a}{2}t+\tfrac{a+b}{2}$.\\
Adaptiv: Simpson$(h)$ vs.\ Simpson$(h/2)$ sch\"atzt lokalen Fehler.

% =================================================
\section{Anfangswertprobleme (ODE)}
% =================================================
$y'=f(t,y),\ y(t_0)=y_0$
\subsection{Euler explizit / implizit}
\fitline{y_{n+1}=y_n+h f(t_n,y_n)\ \ (\text{Ord.\ 1, explizit})}
Stabilit\"at $y'{=}\lambda y$: $|1+h\lambda|\le 1\Rightarrow h\le -2/\lambda$ ($\lambda<0$).\\
\fitline{y_{n+1}=y_n+h f(t_{n+1},y_{n+1})\ \ (\text{implizit, A-stabil})}

\subsection{Heun (Trapez)}
\fitline{\tilde y=y_n+h f(t_n,y_n)}
\fitline{y_{n+1}=y_n+\tfrac{h}{2}\bigl(f(t_n,y_n)+f(t_{n+1},\tilde y)\bigr)}
Ordnung 2.

\subsection{RK4}
\fitline{k_1=f(t_n,y_n)}
\fitline{k_2=f(t_n{+}h/2,\,y_n{+}\tfrac{h}{2}k_1)}
\fitline{k_3=f(t_n{+}h/2,\,y_n{+}\tfrac{h}{2}k_2)}
\fitline{k_4=f(t_n{+}h,\,y_n{+}h k_3)}
\fitline{y_{n+1}=y_n+\tfrac{h}{6}(k_1+2k_2+2k_3+k_4)}
Ordnung 4, 4 Auswertungen/Schritt.

\subsection{Adaptiv \& Steif}
\textbf{Adaptiv}: Dormand-Prince (eingebettetes 4(5)).\\
\textbf{Steif}: gr.\ Spannweite $\max|\lambda|/\min|\lambda|$ $\to$ implizite Verf.\ (BDF, Trapez).

% =================================================
\section{Eigenwerte numerisch}
% =================================================
\textbf{Potenzmethode}: $x_{k+1}=Ax_k/\|Ax_k\|$,\ Rayleigh $\lambda_1\approx x_k\trans Ax_k/x_k\trans x_k$.\\
Rate $|\lambda_2/\lambda_1|$.\\[0.3ex]
\textbf{Inverse Iter.\ mit Shift}: $(A-\mu I)x_{k+1}=x_k$, normieren $\to$ EW nahe $\mu$.\\[0.3ex]
\textbf{QR-Iteration}: $A_0{=}A$;\ $A_k=Q_kR_k,\ A_{k+1}=R_kQ_k$ $\to$ Schur-Form (EW auf Diag).\\
\textbf{Krylov} (Lanczos / Arnoldi) f\"ur grosse d\"unne $A$.

% =================================================
\section{Klausur-Stolperfallen}
% =================================================
\begin{itemize}
\item Ausl\"oschung umformen, nicht direkt.
\item Gauss ohne Pivot bei kleinem $a_{jj}$ $\to$ Pivotsuche.
\item Steifes Problem mit explizitem Euler $\to$ explodiert; impliziter Euler/BDF.
\item Stop-Krit.: \emph{Residuum} besser als \emph{Inkrement}.
\item Zentrale Differenzen schlagen Vorw\"arts (Ordnung und opt.\ $h$).
\item Newton braucht $f'(x_k)\!\ne\!0$ und gute Startn\"ahe.
\item Simpson nur f\"ur gerade Anzahl Teilintervalle.
\item Runge bei \"aquidistanten Knoten $\to$ Tschebyscheff.
\end{itemize}

\end{multicols*}
\end{document}
